\documentclass[11pt,a4paper]{article}
\usepackage[margin=27mm]{geometry}
\usepackage[T1]{fontenc}
\usepackage[utf8]{inputenc}
\usepackage{lmodern,amsmath,amssymb,amsthm,mathtools,booktabs,longtable,array}
\providecommand{\mathscr}[1]{\mathcal{#1}}
\usepackage[dvipsnames]{xcolor}
\usepackage{microtype}
\PassOptionsToPackage{hyphens}{url}
\usepackage[hidelinks]{hyperref}
\Urlmuskip=0mu plus 1mu\relax
\usepackage{enumitem}
\setlist{itemsep=3pt,topsep=5pt}
\setlength{\parskip}{4pt}
\setlength{\parindent}{0pt}
\setlength{\emergencystretch}{3em}
\newtheorem{theorem}{Theorem}[section]
\newtheorem{lemma}[theorem]{Lemma}
\newtheorem{proposition}[theorem]{Proposition}
\newtheorem{corollary}[theorem]{Corollary}
\newtheorem{definition}[theorem]{Definition}
\DeclareMathOperator{\CH}{CH}
\DeclareMathOperator{\Pic}{Pic}
\DeclareMathOperator{\NS}{NS}
\DeclareMathOperator{\End}{End}
\DeclareMathOperator{\Hom}{Hom}
\DeclareMathOperator{\Alb}{Alb}
\providecommand{\tightlist}{\setlength{\itemsep}{0pt}\setlength{\parskip}{0pt}}
\title{JAC1964: primitive incidence polarizations,\newline pencil heights, and the arithmetic complement}
\author{Run JAC1964--20261121\\GPT-6 Astra Pro; collaborative mathematical investigation}
\date{21 November 2026\\Working manuscript: continuously checked findings}
\begin{document}
% Generated from the authored live TSV. Do not edit this file manually.
% Requires longtable, booktabs, array, amsmath; compatible with pdflatex.
\section*{Appendix: finding-to-Grothendieck locator crosswalk}
\addcontentsline{toc}{section}{Finding-to-Grothendieck locator crosswalk}
\begingroup
\small
\setlength{\tabcolsep}{4pt}
\renewcommand{\arraystretch}{1.12}
\setlength{\LTpre}{5pt}
\setlength{\LTpost}{7pt}

This appendix maps every stable finding ID to the supplied archive locus
and, separately, to any actual SGA/EGA/FGA input. The archive keys and
source keys are expanded in full below; they are exact cross-references,
not generic citations. Modern primary references and the precise scope
of each conclusion remain in its proof, so their long bibliographies
are not repeated here.

\textbf{Reading the roles.} A \emph{programme claim} is a repaired or
proved form of the indicated 1964 question. A \emph{supporting lemma}
or \emph{scope obstruction} serves that question. A \emph{modern
application} is explicitly \emph{not a statement in Grothendieck}.
Independently, the final column distinguishes a \emph{direct input}
from a \emph{related framework or convention}; the latter does not
prove the new conclusion. An inherited direct input is identified as
such in the role text. Where none is used, the table says
\emph{no direct SGA/EGA/FGA input}.

Only the supplied cote transcriptions were available. Their scan
numbers are supplied-file locators, not reconstructed manuscript
pagination; no palaeographic verification is claimed. Cote~44 scans
3, 5 and 7 contain unrelated interleaved material. The incidence
quotient is Saito's higher Picard variety up to isogeny; a universal
regular algebraic representative is identified only in the expressly
proved codimension-two cases. No general higher-Picard theorem is
invented or attributed to FGA.

\subsection*{Archive register}
\begin{longtable}{@{}>{\raggedright\arraybackslash}p{.115\linewidth}>{\raggedright\arraybackslash}p{\dimexpr.885\linewidth-2\tabcolsep\relax}@{}}
\toprule
Key & Exact supplied locus and subject \\
\midrule\endfirsthead
\toprule Key & Exact supplied locus and subject (continued) \\
\midrule\endhead
\bottomrule\endfoot
\textsf{A123} & cote 44, scan 2, Th. 1--3 (incidence kernel, dimension bound, quotient and duality); scan 8, letter to Serre of 12 August 1964, first paragraph \\
\textsf{A4} & cote 44, scan 2, Th. 4; scan 8, primitive-decomposition margin; scan 9, normalized family z and divisor \(D=r_*(p^*z\cdot q^*z)\) \\
\textsf{A5} & cote 44, scan 4, Th. 5 (algebraically trivial restrictions and height/intersection formula), including the margin on general fibrations and a local variant \\
\textsf{A6} & cote 44, scan 4, Th. 6: \(2i\leq n-1\) isogeny, \(2i=n\) fixed factor, \(2i\geq n+1\) surjectivity, and the deleted critical finite-kernel clause \\
\textsf{A7} & cote 44, scan 6, Th. 7 (numerical equivalence = \(\tau\)-equivalence; \(\tau\) is not defined on this leaf) \\
\textsf{AODD} & cote 44, scan 8, first paragraph: \(T_\ell(J^i)\) as an odd-cohomology subquotient and \(\dim J^i\leq b_{2i-1}/2\) \\
\textsf{LSTR} & cote 16 no. 7, scan 33, B(4), and scans 34, 36, 38, 40: normalized Lefschetz lowering and primitive strings \\
\textsf{LSEC} & cote 16 no. 7, scan 40, boxed \(\Lambda_Y=\varphi^*\Lambda_X^2\varphi_*\); scan 45, (iv) and (iv bis), algebraic left/right inverses \\
\textsf{LMOT} & cote 16 no. 7, scan 49, algebraic inverse corollary and semisimple-motive paragraph; scan 50, weak-Lefschetz variants \\
\end{longtable}

\subsection*{Exact SGA/EGA source register}
Page numbers are the original printed pagination. No page number has
been guessed where only a section was checked. The three original
volumes are SGA~6, SGA~7~I and SGA~7~II; the volume, exposition and
number in each entry determine the cited statement uniquely.

\begin{longtable}{@{}>{\raggedright\arraybackslash}p{.115\linewidth}>{\raggedright\arraybackslash}p{\dimexpr.885\linewidth-2\tabcolsep\relax}@{}}
\toprule
Key & Exact locus and role \\
\midrule\endfirsthead
\toprule Key & Exact locus and role (continued) \\
\midrule\endhead
\bottomrule\endfoot
\textsf{S6tau} & SGA 6, Exp. XIII, 4.1, printed p. 643, and 4.4, p. 644. Definition/saturation convention for \(\tau\); not a theorem identifying higher-cycle numerical and algebraic equivalence. \\
\textsf{S6div} & SGA 6, Exp. XIII, Theorem 4.6, printed pp. 645--646. Numerically trivial divisors are \(\tau\)-trivial, in every characteristic. \\
\textsf{S6index} & SGA 6, Exp. XIII, Theorem 7.1, printed p. 662. Surface divisor index theorem; it is not higher-codimension or arithmetic Hodge index. \\
\textsf{S7exist} & SGA 7 II, Exp. XVII, 2.2 and Theorem 2.5.1, printed p. 217; characteristic-two qualification on printed p. 219. Lefschetz-pencil definition and existence for degree-d multiples, \(d\geq2\); even ambient dimension avoids the stated degenerate-quadratic exception. \\
\textsf{EGAaxis} & EGA~IV\textsubscript{4}, 17.13.7, as explicitly cited in SGA 7 II, Exp. XVII, 2.2. Transverse axis background only. \\
\textsf{EGAbc} & EGA~III\textsubscript{2}, Theorem 7.7.5(I--II), printed pp. 67--68 (continuous pp. 199--200); Proposition 7.7.10, p. 71 (203); Proposition 7.8.4, pp. 73--74 (205--206). Coherent semicontinuity and base change; on the reduced parameter bases, constant h0 gives local freeness and arbitrary base change, and hence the required section lifting. \\
\textsf{S7cond} & SGA 7 II, Exp. XVIII, 5.3.2, printed p. 287; 5.3.4--5.3.5, p. 288; Proposition 6.2, p. 319. Condition A, extension of direct-image sheaves, and its relation to nontrivial local inertia. \\
\textsf{S7gauss} & SGA 7 II, Exp. XVIII, Corollary 6.4, printed pp. 320--321, and Lemma 6.4.2.3, p. 321. High-degree condition A and the Euler-characteristic/Chern-class calculation. \\
\textsf{S7van} & SGA 7 II, Exp. XVIII, Theorem 6.6, printed pp. 323--325. Spanning and conjugacy of vanishing cycles, with the stated dimension/characteristic hypotheses. \\
\textsf{S7irred} & SGA 7 II, Exp. XVIII, Corollary 6.7.2, printed p. 326; finite-index proof, p. 327. Absolute irreducibility after finite-index monodromy restriction in even ambient dimension. \\
\textsf{S7inv} & SGA 7 II, Exp. XVIII, 6.6.4, referring to 5.6.10. Ambient-image/invariant-cycle framework; this does not by itself algebraize a trace-factorization correspondence. \\
\textsf{S7pl} & SGA 7 II, Exp. XV, 3.3.3, printed p. 191; 3.3.4--3.3.6, p. 192. Vanishing generator up to sign and Picard--Lefschetz coefficient \((-1)^{r+1}\) for relative dimension 2r+1. \\
\textsf{S7biext} & SGA 7 I, Exp. VIII, 4.12.1 and 2.3.11. Symmetric biextension \(\delta(L)=m^*L-\operatorname{pr}_1^*L-\operatorname{pr}_2^*L\) and the anticommutative Cartier-duality comparison. \\
\textsf{S7toric} & SGA 7 I, Exp. IX, 5.5.1--5.5.8, printed pp. 378--379; 9.1--9.2, pp. 418--421; 10.2.3--10.2.8, pp. 440--441; 10.3.5, p. 443; Theorem 10.4, p. 444; 12.5--12.7. Integral all-prime toric monodromy pairing, polarization positivity, ramification factor, and curve normalization. \\
\textsf{S7Griff} & SGA 7 II, Exp. XX, Theorem 4.1, printed p. 355; Remark 4.3, p. 356; \S{}5.1. High-degree hyperplane Griffiths-group obstruction under the theorem's hypotheses; the present application is in characteristic zero. \\
\textsf{S7ret} & SGA 7 II, Exp. XX, 3.3--3.4. Algebraic-retraction condition B and the Lefschetz-lowering mechanism. \\
\end{longtable}

\textbf{Source-check boundary.} The original printed SGA pages recorded
in the source register and live TSV were visually inspected where that
record says so. The section-only entries \textsf{S7inv},
\textsf{S7biext} and \textsf{S7ret}, and the section-only portions of
other entries, were read in a primary-content transcription. The EGA
locator \textsf{EGAaxis} was verified as SGA's explicit reference;
the original EGA theorem was not independently inspected. It is listed
as transverse-axis background, not as a substitute for a proved step.
The accepted 032/Milne inputs are used on the terms stated in the paper.

\subsection*{Complete finding crosswalk}
The inventory contains 105 finding IDs. Rows are combined only when
their archive locators, archive roles, SGA/EGA/FGA locators and
SGA/EGA/FGA roles are identical. A large identical group may be
split across several rows for readability. Each ID is listed
once in this table; the TSV separately records its proof file
and hash. Enumeration and hashes are mechanical coverage checks,
not formal proof verification.

\begin{longtable}{@{}>{\raggedright\arraybackslash}p{.355\linewidth}>{\raggedright\arraybackslash}p{\dimexpr.235\linewidth-2\tabcolsep\relax}>{\raggedright\arraybackslash}p{\dimexpr.410\linewidth-2\tabcolsep\relax}@{}}
\toprule
Finding ID and conclusion & Archive connection & Actual SGA/EGA/FGA input and role \\
\midrule\endfirsthead
\toprule
Finding ID and conclusion & Archive connection & Actual SGA/EGA/FGA input and role (continued) \\
\midrule\endhead
\bottomrule\endfoot
\textbf{JAC-T4.1}. Algebraic primitive projector and nondegenerate incidence summand\par\smallskip \textbf{JAC-T4.2}. Curve-test transfer to J(C)\(\times\)A with exact factor \((g-1)!\) & \textsf{A4}, \textsf{LSTR}.\par \textbf{Supporting lemma.} Technical support for the indicated claim. & \textbf{No direct SGA/EGA/FGA input.} \\
\addlinespace[4pt]
\textbf{JAC-T4.3}. Signed primitive incidence gives a rational polarization & \textsf{A4}.\par \textbf{Programme claim.} Repaired or proved form; the new proof is not attributed to Grothendieck. & \textsf{S7biext}.\par \textbf{Related framework or convention.} Related normalization only; SGA's biextension convention is not the positivity theorem. Positivity is the curve-test deduction. \\
\addlinespace[4pt]
\textbf{JAC-T4.4}. Actual codimension-two regular representative on abelian threefolds & \textsf{A123}, \textsf{A4}.\par \textbf{Programme claim.} Repaired or proved form; the new proof is not attributed to Grothendieck. & \textbf{No direct SGA/EGA/FGA input.} \\
\addlinespace[4pt]
\textbf{JAC-T4.5}. Descent of the primitive incidence polarization & \textsf{A123}, \textsf{A4}.\par \textbf{Supporting lemma.} Technical support for the indicated claim. & \textbf{No direct SGA/EGA/FGA input.} \\
\addlinespace[4pt]
\textbf{JAC-P00}. Rational descent of incidence quotients and Poincar\'{e} families & \textsf{A123}, \textsf{A6}.\par \textbf{Supporting lemma.} Technical support for the indicated claim. & \textbf{No direct SGA/EGA/FGA input.} \\
\addlinespace[4pt]
\textbf{JAC-P01}. Odd-cohomology detection of incidence homomorphisms & \textsf{A123}, \textsf{AODD}.\par \textbf{Programme claim.} Repaired or proved form; the new proof is not attributed to Grothendieck. & \textbf{No direct SGA/EGA/FGA input.} \\
\addlinespace[4pt]
\textbf{JAC-P02}. All three Lefschetz restrictions and the critical fixed trace & \textsf{A6}, \textsf{LSEC}, \textsf{LSTR}.\par \textbf{Programme claim.} Repaired or proved form; the new proof is not attributed to Grothendieck. & \textsf{S7inv}, \textsf{S7ret}.\par \textbf{Related framework or convention.} Related cohomological/algebraic-retraction framework; the trace assertion is proved by rational families plus the blowup formula, not inferred merely from invariant cohomology. \\
\addlinespace[4pt]
\textbf{JAC-P03}. Explicit E\({}^{3}\) test: unsigned incidence is indefinite & \textsf{A4}.\par \textbf{Scope obstruction.} Counterexample or boundary of an unqualified reading. & \textsf{S7biext}.\par \textbf{Related framework or convention.} Sign convention only: P=\(\Delta\)\(-\)axes induces the negative of the usual O(0) polarization. The E\({}^{3}\) counterexample is computed directly. \\
\addlinespace[4pt]
\textbf{JAC-P04}. E\({}^{4}\) hyperplane restrictions expose the homological/algebraic domain gap & \textsf{A5}, \textsf{AODD}.\par \textbf{Scope obstruction.} Counterexample or boundary of an unqualified reading. & \textsf{S7Griff}, \textsf{S7gauss}, \textsf{S7irred}.\par \textbf{Direct input.} Direct application of XX Theorem 4.1/Remark 4.3 with their hypotheses verified; not an all-characteristic Griffiths claim. \\
\addlinespace[4pt]
\textbf{JAC-P05}. CM Fermat counterexample to numerical=\(\tau\) in characteristic zero & \textsf{A7}.\par \textbf{Scope obstruction.} Counterexample or boundary of an unqualified reading. & \textsf{S6tau}, \textsf{S6div}.\par \textbf{Related framework or convention.} XIII supplies the \(\tau\) convention and its divisor boundary only. The higher-cycle counterexample is Bloch's, not a theorem of SGA. \\
\addlinespace[4pt]
\textbf{JAC-P06}. Algebraically trivial cycles over \(\overline{\mathbf F}_p\) are torsion & \textsf{A7}.\par \textbf{Supporting lemma.} Technical support for the indicated claim. & \textsf{S6tau}.\par \textbf{Related framework or convention.} Interpretation of \(\tau\) as saturation only; the torsion argument uses curve Jacobians over finite fields. \\
\addlinespace[4pt]
\textbf{JAC-P07}. Automatic condition A for genuine abelian Lefschetz pencils & \textsf{A5}, \textsf{A6}.\par \textbf{Supporting lemma.} Technical support for the indicated claim. & \textsf{S7cond}, \textsf{S7gauss}, \textsf{S7van}, \textsf{S7irred}.\par \textbf{Direct input.} Direct use of Euler characteristic, spanning/conjugacy and condition A; the Catalan formula and positivity estimate are the deduction. \\
\addlinespace[4pt]
\textbf{JAC-P08}. Algebraic lowering for complete intersections and products & \textsf{LSEC}, \textsf{LMOT}, \textsf{A6}.\par \textbf{Programme claim.} Repaired or proved form; the new proof is not attributed to Grothendieck. & \textbf{No direct SGA/EGA/FGA input.} \\
\addlinespace[4pt]
\textbf{JAC-P09}. Geometric-generic high-degree axis has no variable algebraic classes & \textsf{A5}, \textsf{A6}.\par \textbf{Modern application.} Not a statement in Grothendieck. & \textsf{S7van}.\par \textbf{Related framework or convention.} Vanishing-cycle framework only; the required large orthogonal monodromy is Katz--Pandharipande, not XVIII finite-index symplectic irreducibility. \\
\addlinespace[4pt]
\textbf{JAC-P10}. Exact stable Fermat Hodge-lattice index and Faltings-height gap\par\smallskip \textbf{AR-4}. Elementary analytic central-sign lemma\par\smallskip \textbf{AR-5}. Conditional implication from a positive-normalization height identity & \textsf{A5}.\par \textbf{Modern application.} Not a statement in Grothendieck. & \textbf{No direct SGA/EGA/FGA input.} \\
\addlinespace[4pt]
\textbf{AR-6}. Exact arithmetic slope-method gate\par\smallskip \textbf{AR-8}. Exact Fermat symmetry and Gross--Schoen height formula\par\smallskip \textbf{AR-9}. Positive finite-local Fermat contribution & \textsf{A5}.\par \textbf{Modern application.} Not a statement in Grothendieck. & \textbf{No direct SGA/EGA/FGA input.} \\
\addlinespace[4pt]
\textbf{AR-10}. Actual canonical adelic line representing the primitive height\par\smallskip \textbf{AR-12}. Exact archimedean theta-integral reduction and Fermat period matrix\par\smallskip \textbf{AR-13}. Certified strict positivity of the specified Fermat cycle & \textsf{A5}.\par \textbf{Modern application.} Not a statement in Grothendieck. & \textbf{No direct SGA/EGA/FGA input.} \\
\addlinespace[4pt]
\textbf{AR-18}. Height sign equals an ordering of sum and difference surface heights\par\smallskip \textbf{AR-21}. Every odd rational class comes from one geometrically integral curve\par\smallskip \textbf{AR-21a}. Integral-curve replacement with an even ample correction & \textsf{A5}.\par \textbf{Modern application.} Not a statement in Grothendieck. & \textbf{No direct SGA/EGA/FGA input.} \\
\addlinespace[4pt]
\textbf{AR-21b}. Single-curve Ceresa representation of any odd codimension-two class\par\smallskip \textbf{AR-24}. Every complementary height is a projected Gross--Schoen height\par\smallskip \textbf{AR-27}. The projected-height test can use actual canonical base points & \textsf{A5}.\par \textbf{Modern application.} Not a statement in Grothendieck. & \textbf{No direct SGA/EGA/FGA input.} \\
\addlinespace[4pt]
\textbf{AR-27a}. Original-field quadratic cover with a canonical marked point\par\smallskip \textbf{AR-27b}. Exact projected-diagonal pullback and one-half height factor\par\smallskip \textbf{AR-28}. Exact adelic nef cone for the canonical diagonal metric & \textsf{A5}.\par \textbf{Modern application.} Not a statement in Grothendieck. & \textbf{No direct SGA/EGA/FGA input.} \\
\addlinespace[4pt]
\textbf{AR-29}. Full-Jacobian diagonal decomposition and constraint boundary\par\smallskip \textbf{AR-29a}. Exact elliptic quotient and orthogonal Jacobian complement\par\smallskip \textbf{AR-29b}. Opposite-CM Tate--Picard block and remaining higher-level block & \textsf{A5}.\par \textbf{Modern application.} Not a statement in Grothendieck. & \textbf{No direct SGA/EGA/FGA input.} \\
\addlinespace[4pt]
\textbf{AR-29c}. Canonicalization creates no one-new-Prym term\par\smallskip \textbf{AR-29d}. Scalar-vertical diagonal constraints admit cross negation\par\smallskip \textbf{AR-29e}. Block reflection retains two-A/one-B Picard data & \textsf{A5}.\par \textbf{Modern application.} Not a statement in Grothendieck. & \textbf{No direct SGA/EGA/FGA input.} \\
\addlinespace[4pt]
\textbf{AR-29f}. Canonical one-A/two-B Lefschetz contraction vanishes\par\smallskip \textbf{AR-29g}. Canonical base point does not kill every Picard contraction & \textsf{A5}.\par \textbf{Modern application.} Not a statement in Grothendieck. & \textbf{No direct SGA/EGA/FGA input.} \\
\addlinespace[4pt]
\textbf{JAC-P11}. Positive two-dimensional CM orbit plane\par\smallskip \textbf{JAC-P12}. The tautological ring does not enlarge the primitive grade-one line & \textsf{A5}, \textsf{A7}.\par \textbf{Modern application.} Not a statement in Grothendieck. & \textbf{No direct SGA/EGA/FGA input.} \\
\addlinespace[4pt]
\textbf{JAC-P13}. Canonical actual invariant Chow summand \(N=h^3(J)^G\)\par\smallskip \textbf{JAC-P14}. Full endomorphism algebra acts on N through its determinant & \textsf{AODD}, \textsf{LMOT}, \textsf{A5}.\par \textbf{Modern application.} Not a statement in Grothendieck. & \textbf{No direct SGA/EGA/FGA input.} \\
\addlinespace[4pt]
\textbf{JAC-P15}. Finite-field rational Chow=numerical for geometric abelian type & \textsf{A7}, \textsf{LMOT}.\par \textbf{Programme claim.} Repaired or proved form; the new proof is not attributed to Grothendieck. & \textsf{S6tau}.\par \textbf{Related framework or convention.} \(\tau\) convention only. The finite-field theorem is Kahn plus the accepted Tate input, not SGA XIII's divisor theorem. \\
\addlinespace[4pt]
\textbf{JAC-P16}. B(X) and numerical=homological for geometric abelian-type varieties & \textsf{LMOT}, \textsf{LSEC}, \textsf{A6}.\par \textbf{Programme claim.} Repaired or proved form; the new proof is not attributed to Grothendieck. & \textbf{No direct SGA/EGA/FGA input.} \\
\addlinespace[4pt]
\textbf{JAC-P17}. Actual fields of definition of N, \(\eta\) and T\par\smallskip \textbf{H9-T15}. Conductor-correct CM Selmer rank\par\smallskip \textbf{AR-16}. Scalar-extended Abel--Jacobi injectivity would close the rank gate & \textsf{AODD}, \textsf{A5}.\par \textbf{Modern application.} Not a statement in Grothendieck. & \textbf{No direct SGA/EGA/FGA input.} \\
\addlinespace[4pt]
\textbf{AR-16a}. Exact rational-field injectivity criterion\par\smallskip \textbf{AR-17}. Audit of CM Abel--Jacobi theorems and conductor restrictions\par\smallskip \textbf{AR-19}. Explicit nonzero Katz moment at 5 & \textsf{AODD}, \textsf{A5}.\par \textbf{Modern application.} Not a statement in Grothendieck. & \textbf{No direct SGA/EGA/FGA input.} \\
\addlinespace[4pt]
\textbf{AR-23}. The actual Fermat cycle spans the proved Selmer target & \textsf{AODD}, \textsf{A5}.\par \textbf{Modern application.} Not a statement in Grothendieck. & \textbf{No direct SGA/EGA/FGA input.} \\
\addlinespace[4pt]
\textbf{JAC-P18}. A curve-times-J fourfold obstruction to algebraic-radical vanishing & \textsf{A123}, \textsf{AODD}, \textsf{A5}.\par \textbf{Modern application.} Not a statement in Grothendieck. & \textbf{No direct SGA/EGA/FGA input.} \\
\addlinespace[4pt]
\textbf{JAC-P19}. An ample complete-intersection axis has no abelian rational domination & \textsf{A5}.\par \textbf{Scope obstruction.} Counterexample or boundary of an unqualified reading. & \textbf{No direct SGA/EGA/FGA input.} \\
\addlinespace[4pt]
\textbf{JAC-P20}. Actual total-cohomology injection from an abelian-type Chow splitting & \textsf{A5}, \textsf{A6}, \textsf{LMOT}.\par \textbf{Modern application.} Not a statement in Grothendieck. & \textbf{No direct SGA/EGA/FGA input.} \\
\addlinespace[4pt]
\textbf{H9-C1}. Ordinary quadric pencil refutes the uncorrected height identification & \textsf{A5}.\par \textbf{Scope obstruction.} Counterexample or boundary of an unqualified reading. & \textsf{S7exist}, \textsf{S7cond}.\par \textbf{Direct input.} The pencil satisfies the precise Lefschetz definition but has zero vanishing cohomology; geometrically integral fibres alone do not remove higher vertical cycles. \\
\addlinespace[4pt]
\textbf{H9-T1}. Corrected height formula with a vertical intersection matrix & \textsf{A5}.\par \textbf{Programme claim.} Repaired or proved form; the new proof is not attributed to Grothendieck. & \textbf{No direct SGA/EGA/FGA input.} \\
\addlinespace[4pt]
\textbf{H9-T2}. Specialization injectivity removes local height corrections & \textsf{A5}.\par \textbf{Programme claim.} Repaired or proved form; the new proof is not attributed to Grothendieck. & \textsf{S7cond}, \textsf{S7pl}.\par \textbf{Direct input.} Direct local-cohomology input; cap-product admissibility and the height identity are proved in the fragment. \\
\addlinespace[4pt]
\textbf{H9-T3}. Exact quotient/radical statement and the domain correction & \textsf{A123}, \textsf{A5}.\par \textbf{Scope obstruction.} Counterexample or boundary of an unqualified reading. & \textsf{S7Griff}.\par \textbf{Direct input.} XX Theorem 4.1/Remark 4.3 supports the explicit domain obstruction; the positive-quotient implication is linear algebra plus N\'{e}ron--Tate positivity. \\
\addlinespace[4pt]
\textbf{H9-T4}. Polarization of the fixed primitive part of a smooth ample section & \textsf{A4}, \textsf{A6}, \textsf{LSEC}.\par \textbf{Programme claim.} Repaired or proved form; the new proof is not attributed to Grothendieck. & \textsf{S7inv}.\par \textbf{Related framework or convention.} Fixed/variable cohomological framework; the sign is transferred to J(C)\(\times\)X and uses accepted 032, not an HSC assumption for Y. \\
\addlinespace[4pt]
\textbf{H9-T5}. One out-of-range Newton slope kills the variable incidence factor & \textsf{A4}, \textsf{AODD}.\par \textbf{Supporting lemma.} Technical support for the indicated claim. & \textsf{S7irred}.\par \textbf{Direct input.} Monodromy hypothesis has this exact standard-pencil locus. The Newton exclusion is a separate Frobenius/good-reduction argument. \\
\addlinespace[4pt]
\textbf{H9-L1}. A rational scalar locally square at almost every prime is positive & \textsf{A4}.\par \textbf{Supporting lemma.} Technical support for the indicated claim. & \textbf{No direct SGA/EGA/FGA input.} \\
\addlinespace[4pt]
\textbf{H9-L2}. Positive rational rank-one toric monodromy coefficient & \textsf{A4}, \textsf{A5}.\par \textbf{Supporting lemma.} Technical support for the indicated claim. & \textsf{S7toric}, \textsf{S7biext}.\par \textbf{Direct input.} Direct derivation from the integral pairing, the precise polarization direction and the curve normalization; \(q_0=u/s=(us)/s^2\). \\
\addlinespace[4pt]
\textbf{H9-T6}. Variable incidence polarization from all-prime monodromy & \textsf{A4}.\par \textbf{Programme claim.} Repaired or proved form; the new proof is not attributed to Grothendieck. & \textsf{S7pl}, \textsf{S7irred}, \textsf{S7toric}, \textsf{S7biext}.\par \textbf{Direct input.} Direct all-prime comparison \(a_\ell^2=q_0/(ec)\) plus H9-L1 forces the single rational incidence scalar \(c>0\). \\
\addlinespace[4pt]
\textbf{H9-T6a}. Concrete generic-pencil primitive polarization theorem & \textsf{A4}, \textsf{A6}, \textsf{LSEC}.\par \textbf{Programme claim.} Repaired or proved form; the new proof is not attributed to Grothendieck. & \textsf{S7exist}, \textsf{S7van}, \textsf{S7irred}, \textsf{S7pl}, \textsf{S7toric}.\par \textbf{Direct input.} Direct verification of the hypotheses of H9-T6; \(d\geq2\) multiples supply genuine pencils; a separate condition A assumption is unnecessary when the variable incidence factor is nonzero. \\
\addlinespace[4pt]
\textbf{H9-T7}. Whole original-pencil height quotient equals ambient plus axis forms & \textsf{A5}, \textsf{A6}, \textsf{LSEC}.\par \textbf{Modern application.} Not a statement in Grothendieck. & \textsf{S7cond}.\par \textbf{Direct input.} Condition A controls all vertical classes and admissibility, automatically available for the abelian pencil by P07. \\
\addlinespace[4pt]
\textbf{H9-T8}. Complete all-homological height positivity for abelian-fourfold pencils & \textsf{A5}.\par \textbf{Modern application.} Not a statement in Grothendieck. & \textsf{S6index}, \textsf{S6div}, \textsf{S7cond}.\par \textbf{Direct input.} Direct all-characteristic surface index theorem and numerical=homological for divisors, combined with H9-T7. \\
\addlinespace[4pt]
\textbf{H9-T8a}. Arbitrary original-pencil axes in characteristic zero & \textsf{A5}.\par \textbf{Modern application.} Not a statement in Grothendieck. & \textsf{S7cond}.\par \textbf{Direct input.} Only the admissibility input is SGA. Positivity of the axis form is classical Hodge--Riemann, not a new SGA consequence. \\
\addlinespace[4pt]
\textbf{H9-T8b}. Exact reduction of the fixed-axis sign to finite fields & \textsf{A5}, \textsf{LSTR}.\par \textbf{Modern application.} Not a statement in Grothendieck. & \textbf{No direct SGA/EGA/FGA input.} Smooth proper base change and cycle-class compatibility are used through the exact primary Milne exposition listed; no unverified original SGA locator or reverse algebraic specialization is asserted. \\
\addlinespace[4pt]
\textbf{H9-T9}. Incidence height descends exactly to the Lang--N\'{e}ron quotient & \textsf{A123}, \textsf{A4}, \textsf{A5}, \textsf{A6}.\par \textbf{Programme claim.} Repaired or proved form; the new proof is not attributed to Grothendieck. & \textbf{No direct SGA/EGA/FGA input.} No direct SGA result replaces the height comparison. The primitive-polarization input may be supplied upstream by H9-T6/T6a. \\
\addlinespace[4pt]
\textbf{H9-T10}. Actual Ab\({}^{2}\) for generic threefold sections & \textsf{A123}, \textsf{A4}.\par \textbf{Programme claim.} Repaired or proved form; the new proof is not attributed to Grothendieck. & \textsf{S7pl}, \textsf{S7irred}, \textsf{S7toric}.\par \textbf{Direct input.} Inherited directly through the repeated variable-factor monodromy proof; no unproved all-codimension representative is substituted. \\
\addlinespace[4pt]
\textbf{H9-T11}. Elliptic-product reduction with an exact numerical kernel & \textsf{A4}, \textsf{LMOT}.\par \textbf{Programme claim.} Repaired or proved form; the new proof is not attributed to Grothendieck. & \textsf{S7biext}.\par \textbf{Related framework or convention.} Curve sign convention only; the separate D(X) requirement is not inferred from incidence positivity. \\
\addlinespace[4pt]
\textbf{H9-T12}. Finite-base-change norm decomposition and preserved primitive trace & \textsf{A5}, \textsf{A6}.\par \textbf{Modern application.} Not a statement in Grothendieck. & \textsf{S7cond}.\par \textbf{Direct input.} Condition A controls admissibility of those separable covers. No full-hom positivity after every finite cover is claimed; T17 refutes it. \\
\addlinespace[4pt]
\textbf{H9-T13}. Primitive polarization for generic odd complete intersections & \textsf{A4}, \textsf{A6}, \textsf{LSEC}.\par \textbf{Programme claim.} Repaired or proved form; the new proof is not attributed to Grothendieck. & \textsf{S7exist}, \textsf{S7van}, \textsf{S7irred}, \textsf{S7pl}, \textsf{S7toric}.\par \textbf{Direct input.} Terminal-pencil monodromy plus fixed-part transfer with factor \((g+c-1)!/c!\); no application to the nonample product-axis construction. \\
\addlinespace[4pt]
\textbf{H9-T14}. Constant finite-field varieties after arbitrary extension & \textsf{A7}, \textsf{A5}, \textsf{A6}.\par \textbf{Modern application.} Not a statement in Grothendieck. & \textsf{S6tau}.\par \textbf{Related framework or convention.} \(\tau\) interpretation only. Rational Chow=numerical is not claimed over arbitrary transcendental extensions; algebraic equivalence is the conclusion. \\
\addlinespace[4pt]
\textbf{H9-T16}. Semistable axis degeneration and precise nonlifting obstruction & \textsf{A4}, \textsf{A5}, \textsf{LSEC}.\par \textbf{Scope obstruction.} Counterexample or boundary of an unqualified reading. & \textbf{No direct SGA/EGA/FGA input.} No exact SGA/EGA/FGA reverse-Chow-specialization theorem is asserted; the relevant weight-complex inputs are the modern sources listed. \\
\addlinespace[4pt]
\textbf{H9-T17}. Finite covers can create an indefinite norm-zero height plane & \textsf{A5}.\par \textbf{Scope obstruction.} Counterexample or boundary of an unqualified reading. & \textsf{S7cond}, \textsf{S7exist}.\par \textbf{Direct input.} SGA supplies genuine pencils and admissibility after covers unramified above critical values, not a positivity theorem for every finite cover. \\
\addlinespace[4pt]
\textbf{H9-T18}. Finite-cover primitive positivity is exactly a variable fourfold form & \textsf{A4}, \textsf{A5}, \textsf{LSEC}.\par \textbf{Modern application.} Not a statement in Grothendieck. & \textsf{S6index}, \textsf{S6div}, \textsf{S7cond}.\par \textbf{Direct input.} Divisor index and condition A enter the correction/primitive argument; no HSC for the variable incidence-fourfold block is inferred from its abelian-type ambient space. \\
\addlinespace[4pt]
\textbf{H9-T19}. Geometric-generic tame cyclic rational covers add no height directions & \textsf{A5}, \textsf{A6}.\par \textbf{Modern application.} Not a statement in Grothendieck. & \textsf{S7cond}, \textsf{S7irred}, \textsf{S7pl}.\par \textbf{Direct input.} Condition A and finite-index irreducibility are direct inputs. New moving-branch irreducibility uses Katz middle convolution, not an invented SGA orthogonal-monodromy theorem. \\
\addlinespace[4pt]
\textbf{H9-T19a}. Full geometric-generic degree-N rational-map cover extension & \textsf{A5}, \textsf{A6}.\par \textbf{Modern application.} Not a statement in Grothendieck. & \textsf{S7cond}, \textsf{S7irred}, \textsf{S7pl}.\par \textbf{Direct input.} Inherited direct SGA inputs through T19; the passage from cyclic subfamily to generic rational maps is a separate projector/monodromy deduction. \\
\addlinespace[4pt]
\textbf{H9-T20}. Generic positive-genus Kummer covers have the exact height norm formula & \textsf{A5}, \textsf{A6}.\par \textbf{Modern application.} Not a statement in Grothendieck. & \textsf{S7cond}, \textsf{S7irred}, \textsf{S7pl}.\par \textbf{Direct input.} Inherited condition-A and monodromy inputs from T19. Both fixed and variable moving Kummer pieces require the separately cited Katz middle-convolution argument; no result for every specialized cover is asserted. \\
\addlinespace[4pt]
\textbf{H9-T21}. Wild rational scaling and every generic rational-cover degree & \textsf{A5}, \textsf{A6}.\par \textbf{Modern application.} Not a statement in Grothendieck. & \textsf{S7cond}, \textsf{S7irred}, \textsf{S7pl}.\par \textbf{Direct input.} SGA supplies the old pencil's condition A, transvections and finite-index irreducibility. Wild moving-parameter monodromy is the Fourier-transform theorem of Katz listed separately. \\
\addlinespace[4pt]
\textbf{H9-T22}. Wild simple-pole covers and every generic hyperelliptic genus & \textsf{A5}, \textsf{A6}.\par \textbf{Modern application.} Not a statement in Grothendieck. & \textsf{S7cond}, \textsf{S7irred}, \textsf{S7pl}.\par \textbf{Direct input.} Old-pencil SGA inputs are inherited through T21; Fourier-monodromy, the actual wild normalization, and the full double-cover parameter argument are distinct steps, not an attributed SGA theorem. \\
\addlinespace[4pt]
\textbf{H9-T23}. Cyclic triple-cover families of every genus at least two\par\smallskip \textbf{H9-T23a}. Composite positive-genus cover families in degrees 2M and 3M & \textsf{A5}, \textsf{A6}.\par \textbf{Modern application.} Not a statement in Grothendieck. & \textsf{S7cond}, \textsf{S7irred}, \textsf{S7pl}.\par \textbf{Direct input.} Inherited condition A, transvections and finite-index irreducibility concern the original pencil. Moving Kummer or Artin--Schreier monodromy is supplied by Katz, not by a new SGA assertion. \\
\addlinespace[4pt]
\textbf{H9-T24}. Full generic triple-cover space in every genus at least two & \textsf{A5}, \textsf{A6}.\par \textbf{Modern application.} Not a statement in Grothendieck. & \textsf{S7cond}, \textsf{S7irred}, \textsf{S7pl}, \textsf{EGAbc}.\par \textbf{Direct input.} Original-pencil condition A and finite-index monodromy are inherited through T23. Direct EGA coherent base change supplies local freeness and the lift of a cubic section along the balancing extension. \\
\addlinespace[4pt]
\textbf{H9-T25}. Full generic covers above the nonspecial degree bound & \textsf{A5}, \textsf{A6}.\par \textbf{Modern application.} Not a statement in Grothendieck. & \textsf{S7cond}, \textsf{S7irred}, \textsf{S7pl}, \textsf{EGAbc}.\par \textbf{Direct input.} The original-pencil monodromy is inherited through T23a. Direct EGA coherent base change makes the two-section parameter a vector bundle and permits its universal construction. \\
\addlinespace[4pt]
\textbf{H9-T26}. Full generic elliptic triple covers by Radon transform & \textsf{A5}, \textsf{A6}.\par \textbf{Modern application.} Not a statement in Grothendieck. & \textsf{S7cond}, \textsf{S7irred}, \textsf{S7pl}, \textsf{S7exist}.\par \textbf{Direct input.} The original pencil uses condition A and finite-index irreducibility. The ordinary-node formula is tensored with the lisse coefficient stalk; cubic-pencil existence and monodromy handle the fixed ambient block. Radon simplicity is a separate modern input. \\
\addlinespace[4pt]
\textbf{AR-1}. Primitive Beauville-grade reduction of the height question & \textsf{A4}, \textsf{A5}, \textsf{LSTR}.\par \textbf{Modern application.} Not a statement in Grothendieck. & \textbf{No direct SGA/EGA/FGA input.} \\
\addlinespace[4pt]
\textbf{AR-1a}. No extra lower-grade radical on an abelian threefold & \textsf{A5}, \textsf{LSTR}.\par \textbf{Modern application.} Not a statement in Grothendieck. & \textbf{No direct SGA/EGA/FGA input.} \\
\addlinespace[4pt]
\textbf{AR-2}. Exact orthogonal Picard Chow summand and remaining complementary gate & \textsf{A123}, \textsf{A4}, \textsf{A5}, \textsf{LMOT}.\par \textbf{Modern application.} Not a statement in Grothendieck. & \textbf{No direct SGA/EGA/FGA input.} \\
\addlinespace[4pt]
\textbf{AR-3}. Explicit CM group-algebra complementary motive & \textsf{AODD}, \textsf{A5}, \textsf{LMOT}.\par \textbf{Modern application.} Not a statement in Grothendieck. & \textbf{No direct SGA/EGA/FGA input.} \\
\addlinespace[4pt]
\textbf{AR-7}. A nonzero Fermat class in the complementary motive & \textsf{A7}, \textsf{AODD}, \textsf{A5}.\par \textbf{Modern application.} Not a statement in Grothendieck. & \textsf{S6tau}.\par \textbf{Related framework or convention.} \(\tau\) convention only; the nonzero Griffiths class is Bloch's theorem. \\
\addlinespace[4pt]
\textbf{AR-11}. Positive central derivative of the actual Fermat L-function & \textsf{A5}, \textsf{AODD}.\par \textbf{Modern application.} Not a statement in Grothendieck. & \textbf{No direct SGA/EGA/FGA input.} \\
\addlinespace[4pt]
\textbf{AR-14}. Positive arithmetic subspace containing a Griffiths class & \textsf{A123}, \textsf{A4}, \textsf{A5}.\par \textbf{Modern application.} Not a statement in Grothendieck. & \textbf{No direct SGA/EGA/FGA input.} \\
\addlinespace[4pt]
\textbf{AR-15}. Algebraically trivial complementary cycles lie in the entire height radical & \textsf{A123}, \textsf{A5}, \textsf{A7}, \textsf{AODD}.\par \textbf{Modern application.} Not a statement in Grothendieck. & \textbf{No direct SGA/EGA/FGA input.} \\
\addlinespace[4pt]
\textbf{AR-15a}. Grade-zero parameter-Jacobian obstruction for the algebraic radical & \textsf{A123}, \textsf{AODD}, \textsf{LMOT}.\par \textbf{Modern application.} Not a statement in Grothendieck. & \textbf{No direct SGA/EGA/FGA input.} \\
\addlinespace[4pt]
\textbf{AR-20}. Finite-specialization route has an exact obstruction & \textsf{A7}, \textsf{A5}.\par \textbf{Modern application.} Not a statement in Grothendieck. & \textsf{S6tau}.\par \textbf{Related framework or convention.} \(\tau\) convention only; the specialization vanishing uses modern Kahn plus accepted full Tate, not SGA XIII's divisor theorem. \\
\addlinespace[4pt]
\textbf{AR-20a}. Every good finite-field specialization of N has zero rational Chow group & \textsf{A7}, \textsf{AODD}.\par \textbf{Modern application.} Not a statement in Grothendieck. & \textsf{S6tau}.\par \textbf{Related framework or convention.} Convention only; no direct SGA finite-field Chow theorem is claimed. \\
\addlinespace[4pt]
\textbf{AR-20b}. A positive global class is invisible in every good finite specialization & \textsf{A7}, \textsf{A5}.\par \textbf{Modern application.} Not a statement in Grothendieck. & \textbf{No direct SGA/EGA/FGA input.} \\
\addlinespace[4pt]
\textbf{AR-22}. Divisor Hodge-index bounds leave the desired cubic sign unoriented & \textsf{A5}.\par \textbf{Modern application.} Not a statement in Grothendieck. & \textsf{S6index}.\par \textbf{Direct input.} Classical geometric surface index supplies \(t\geq0\) only; the arithmetic inequality is Yuan--Zhang, not SGA. \\
\addlinespace[4pt]
\textbf{AR-25}. Exact Fermat tensor-diagonal and sixteen-intersection reduction\par\smallskip \textbf{AR-26}. The full Hermitian divisor-HI system has a sign symmetry & \textsf{A5}, \textsf{AODD}, \textsf{LMOT}.\par \textbf{Modern application.} Not a statement in Grothendieck. & \textbf{No direct SGA/EGA/FGA input.} \\
\addlinespace[4pt]
\end{longtable}
\endgroup

\end{document}
